% !TeX root = ../../book.tex
\begin{chapex}
\label{cqTranslatePropositionalFormulaeToEnglish}
对于给定的 $n \in \mathbb{N}$，设 $p$ 表示命题`$n$是偶数'，设 $q$ 表示命题`$n$是质数'，设 $r$ 表示命题 `$n = 2$'。 对于以下每个命题公式，将其翻译成简单的中文并确定是否对所有 $n \in \mathbb{N}$ 都为真，还是对于 $n$ 的某些值为真，某些值为假，抑或是对于所有 $n \in \mathbb{N}$ 都为假。

\begin{enumerate}[(a)]
\item $(p \wedge q) \Rightarrow r$
\item $q \wedge (\neg r) \Rightarrow (\neg p)$
\item $((\neg p) \vee (\neg q)) \vee (\neg r)$
\item $(p \wedge q) \wedge (\neg r)$
\end{enumerate}
\end{chapex}

\begin{chapex}
对于以下每个中文陈述，将其翻译成符号命题公式。公式中的命题变量应尽可能代表最简单的命题。
\begin{enumerate}[(a)]
\item 豚鼠很安静，但饥饿时会大声吵叫。
\item $2$ 是偶数并不影响它仍是一个质数。
\item $\sqrt{2}$ 不可能是整数，因为它是无理数。
\end{enumerate}
\end{chapex}

\begin{chapex}
设 $p$ 和$q$ 为命题，并假设 $p \Rightarrow (\neg q)$ 为真，而 $(\neg q) \Rightarrow p$ 为假。以下哪些是正确的，哪些是错误的？
\begin{enumerate}[(a)]
\item $q$ 为假是 $p$ 为真的必要条件。
\item $q$ 为假是 $p$ 为真的充分条件。
\item $p$ 为真是 $q$ 为假的必要条件。
\item $p$ 为真是 $p$ 为假的充分条件。
\end{enumerate}
\end{chapex}

在\Crefrange{cqStructureProofBegin}{cqStructureProofEnd}中，使用\Cref{secPropositionalLogic}中逻辑运算符的定义来描述应遵循哪些步骤来证明问题中的命题公式；字母 $p$、$q$、$r$ 和 $s$ 是命题变量。

\begin{chapex}
\label{cqStructureProofBegin}
$(p \wedge q) \Rightarrow (\neg r)$
\end{chapex}

\begin{chapex}
$(p \vee q) \Rightarrow (r \Rightarrow s)$
\end{chapex}

\begin{chapex}
$(p \Rightarrow q) \Leftrightarrow (\neg p \Rightarrow \neg q)$
\end{chapex}

\begin{chapex}
\label{cqStructureProofEnd}
$(p \wedge (\neg q)) \vee (q \wedge (\neg p))$
\end{chapex}

% 变量和量词 %

\begin{chapex}
找到一个完全不涉及任何变量的简单中文陈述，相当于逻辑公式
$\forall a \in \mathbb{Q},\, \forall b \in \mathbb{Q},\, (a < b \Rightarrow \exists c \in \mathbb{R},\, [a < c < b ~\wedge~ \neg (c \in \mathbb{Q})])$。然后以逻辑公式的结构指导该命题的证明。
\end{chapex}

\begin{chapex}
找出一个与下面的命题等价的纯符号逻辑公式，并证明：
``\textit{无论你选择哪个整数，都会有一个比它大的整数。}''
\end{chapex}

\begin{chapex}
设 $X$ 为集合，$p(x)$ 为谓词。找到一个逻辑公式表示`恰好有两个元素 $x \in X$ 使得 $p(x)$ 为真'。使用这个逻辑公式的结构来描述应该如何构造一个证明，并使用这个结构来证明恰好有两个实数 $x$ 使得 $x^2=1$。
\end{chapex}

% 逻辑等价 %

\begin{chapex}
证明
\[
p \Leftrightarrow q \equiv (p \Rightarrow q) \wedge ((\neg p) \Rightarrow (\neg q))
\]
这种逻辑等价如何帮助您证明`$p$ 当且仅当 $q$'形式的陈述？
\end{chapex}

\begin{chapex}
使用真值表证明 $p \Rightarrow q \not\equiv q \Rightarrow p$。给出一个命题 $p$ 和 $q$ 的例子，使得 $p \Rightarrow q$ 为真，但 $q \Rightarrow p$ 为假。
\end{chapex}

在\Crefrange{cqTruthTablesBegin}{cqTruthTablesEnd}中，找到一个逻辑公式，其真值表中的列如图所示。

\begin{chapex}
\label{cqTruthTablesBegin}
\begin{tabular}{cc|c}
$p$ & $q$ & \hspace{50pt} \\ \hline
\TT & \TT & \FF \\
\TT & \FF & \TT \\
\FF & \TT & \TT \\
\FF & \FF & \FF \\
\end{tabular}
\end{chapex}

\begin{chapex}
\begin{tabular}{cc|c}
$p$ & $q$ & \hspace{50pt} \\ \hline
\TT & \TT & \TT \\
\TT & \FF & \FF \\
\FF & \TT & \TT \\
\FF & \FF & \FF \\
\end{tabular}
\end{chapex}

\begin{chapex}
\begin{tabular}{ccc|c}
$p$ & $q$ & $r$ & \hspace{50pt} \\ \hline
\TT & \TT & \TT & \TT \\
\TT & \TT & \FF & \TT \\
\TT & \FF & \TT & \FF \\
\TT & \FF & \FF & \FF \\
\FF & \TT & \TT & \FF \\
\FF & \TT & \FF & \FF \\
\FF & \FF & \TT & \TT \\
\FF & \FF & \FF & \TT \\
\end{tabular}
\end{chapex}

\begin{chapex}
\label{cqTruthTablesEnd}
\begin{tabular}{ccc|c}
$p$ & $q$ & $r$ & \hspace{50pt} \\ \hline
\TT & \TT & \TT & \TT \\
\TT & \TT & \FF & \FF \\
\TT & \FF & \TT & \FF \\
\TT & \FF & \FF & \FF \\
\FF & \TT & \TT & \TT \\
\FF & \TT & \FF & \FF \\
\FF & \FF & \TT & \TT \\
\FF & \FF & \FF & \FF \\
\end{tabular}
\end{chapex}

\begin{chapex}
新的逻辑运算符 $\uparrow$ 由以下规则定义：
\begin{enumerate}[(i)]
\item 如果从 $p$ 为真的假设中可以引出矛盾，则 $p \uparrow q$ 为真；
\item 如果从 $q$ 为真的假设中可以引出矛盾，则 $p \uparrow q$ 为真；
\item 如果 $r$ 是任意命题，并且如果 $p \uparrow q$、$p$ 和 $q$ 都为真，则 $r$ 为真。
\end{enumerate}

这个问题探讨了这个奇怪的新逻辑运算符。
\begin{enumerate}[(a)]
\item 证明 $p \uparrow p \equiv \neg p$，并推导出 $((p \uparrow p) \uparrow (p \uparrow p)) \equiv p$。
\item 证明 $p \vee q \equiv (p \uparrow p) \uparrow (q \uparrow q)$ 和 $p \wedge q \equiv (p \uparrow q) \uparrow (p \uparrow q)$。
\item 仅使用逻辑运算符 $\uparrow$ 构建等价于 $p \Rightarrow q$ 的命题公式。
\end{enumerate}
\end{chapex}

\begin{chapex}
设 $X$ 为 $\mathbb{Z}$ 或 $\mathbb{Q}$，并定义逻辑公式 $p$ 为:
$$\forall x \in X,\, \exists y \in X,\, (x<y \wedge [\forall z \in X,\, \neg (x<z \wedge z<y)])$$
将 $\neg p$ 写为最大否定逻辑公式。证明当 $X = \mathbb{Z}$ 时 $p$ 为真，当 $X = \mathbb{Q}$ 时 $p$ 为假。
\end{chapex}

\begin{chapex}
使用 \Cref{defUniqueExistentialQuantifier} 写出一个最大否定逻辑公式，等价于 $\neg \exists! x \in X,\, p(x)$。描述该等式给出的用于证明不存在唯一的 $x \in X$ 使得 $p(x)$ 为真的策略，并使用该策略来证明，对于所有 $a \in \mathbb{R }$，如果 $a \ne -1$ 则不存在唯一的 $x \in \mathbb{R}$，使得 $x^4-2ax^2+a^2-1=0$。
\end{chapex}

\begin{chapex}
定义一个新的量词 $\forall !$，使得德摩根量词定律 (\Cref{thmDeMorganQuantifiers}) 成立，并将 $\forall$ 和 $\exists$ 分别替换为 $\forall!$ 和 $\exists!$。
\end{chapex}

\subsection*{判断题}

\tfquestiontext{cqLogicTFBegin}{cqLogicTFEnd}

\begin{chapex} % True
\label{cqLogicTFBegin}
每个蕴涵式都逻辑等价于它的逆否命题。
\end{chapex}

\begin{chapex} % False
每个蕴涵式都逻辑等价于它的逆命题。
\end{chapex}

\begin{chapex} % True
每个逻辑运算符仅为合取和否定的命题公式逻辑等价于逻辑运算符仅为析取和否定的命题公式。
\end{chapex}

\begin{chapex} % False
每个逻辑运算符仅为合取的命题公式逻辑等价于逻辑运算符仅为析取的命题公式。
\end{chapex}

\begin{chapex} % False
公式 $p \wedge (q \vee r)$ 与 $(p \wedge q) \vee r$ 逻辑等价。
\end{chapex}

\begin{chapex} % True
公式 $p \vee (q \vee r)$ 与 $(p \vee q) \wedge (p \vee r)$ 逻辑等价。
\end{chapex}

\begin{chapex} % False
逻辑公式 $\neg \forall x \in X,\, \forall y \in Y,\, p(x,y)$ 与 $\exists x \in X,\, \forall y \in Y,\ , p(x,y)$ 逻辑等价。
\end{chapex}

\begin{chapex} % False
$\neg \forall x \ge 0,\, \exists y \in \mathbb{R},\, y^2=x$ 逻辑等价于 $\forall x < 0,\, \exists y \not\in \mathbb{R},\, y^2 \ne x$。
\end{chapex}

\begin{chapex} % True
$\neg \forall x \ge 0,\, \exists y \in \mathbb{R},\, y^2=x$ 逻辑等价于 $\exists x \ge 0,\, \forall y \in \mathbb{R},\, y^2 \ne x$.
\end{chapex}

\begin{chapex} % False
\label{cqLogicTFEnd}
$\neg \forall x \ge 0,\, \exists y \in \mathbb{R},\, y^2=x$ 逻辑等价于 $\exists x < 0,\, \forall y \in \mathbb{R},\, y^2 = x$.
\end{chapex}

\subsection*{可能性问题}

\asnquestiontext{cqLogicASNBegin}{cqLogicASNEnd}

\begin{chapex} % Sometimes
\label{cqLogicASNBegin}
设 $p$ 和 $q$ 为命题并假设 $p$ 为真。那么 $p \Rightarrow q$ 为真。
\end{chapex}

\begin{chapex} % Always
设 $p$ 和 $q$ 为命题并假设 $p$ 为假。那么 $p \Rightarrow q$ 为真。
\end{chapex}

\begin{chapex} % Always
设 $X$ 和 $Y$ 为集合，并设 $p(x,y)$ 是具有自由变量 $x \in X$ 和 $y \in Y$ 的谓词。那么逻辑公式 $\forall x \in X,\, \forall y \in Y,\, p(x,y)$ 与 $\forall y \in Y,\, \forall x \in X,\, p(x,y)$ 逻辑等价。
\end{chapex}

\begin{chapex} % Sometimes
设 $X$ 和 $Y$ 为集合，并设 $p(x,y)$ 是具有自由变量 $x \in X$ 和 $y \in Y$ 的谓词。那么逻辑公式 $\forall x \in X,\, \exists y \in Y,\, p(x,y)$ 与 $\exists y \in Y,\, \forall x \in X,\, p(x,y)$ 逻辑等价。
\end{chapex}

\begin{chapex} % Never
\label{cqLogicASNEnd}
设 $X$ 和 $Y$ 为集合，并设 $p(x,y)$ 是具有自由变量 $x \in X$ 和 $y \in Y$ 的谓词。那么逻辑公式 $\forall x \in X,\, \forall y \in Y,\, p(x,y)$ 与 $\exists x \in X,\, \exists y \in Y,\, \neg p(x,y)$ 逻辑等价。
\end{chapex}